\section{Desire and the Gift of God}

The heart learns something of its capacity for God through the goods it receives. Friendship frees us from loneliness, and understanding a truth opens possibilities we had not seen. These real goods can awaken a desire that reaches beyond them. We want the person we love to live, and we want to understand more fully what we have begun to know. We can receive these goods gratefully while discovering how much still exceeds our grasp.

Aquinas asks whether any created good can constitute perfect human happiness. His answer turns upon the range of our desire: a finite good offers a real perfection but cannot contain the fullness of good toward which the will reaches (ST I--II, q.~2, a.~8). Wealth, honor, power, and even the whole created universe remain finite. However much we receive from them, our desire still reaches beyond what they can give. Its completion lies in the source from whom they receive their goodness. We can therefore give thanks for a created good while acknowledging that it leaves something to be desired.

We often expect a created good to satisfy that further desire. We may ask another person's affection to secure us against every fear. When affection cannot do this, we begin to blame the person who loves us or to demand ever more assurance. A gift that should bring joy becomes an occasion of anxiety for both people. Hope in God releases the other person from this demand for absolute security. We can attend more freely to his own good when we are less occupied with asking him to answer every need of ours.

The intellect has its own movement of desire. Aquinas describes the wonder that arises when we know an effect but have not reached its cause. An answer can produce further inquiry because knowing that something happens leaves us wanting to understand why. He brings this movement to its highest fulfillment in the vision of God (ST I--II, q.~3, a.~8). While our knowledge stops at created effects, we still desire to know their first cause. Aquinas therefore places perfect happiness in knowing God himself. The desire to understand an effect is finally satisfied when its divine cause is directly known.

Yet desire does not give us a power proportionate to its promised fulfillment. Aquinas teaches that no created intellect, human or angelic, can see the divine essence by its own natural strength (ST I, q.~12, a.~4). We cannot make ourselves capable of heaven by becoming more clever, more refined, or more intense in our longing. The supernatural end is freely given, and God supplies the grace by which it is attained. Our dependence upon grace also gives grounds for confidence. A person whose earthly abilities are modest is not thereby placed farther from the heart of God than a person of extraordinary talent.

Grace makes a saving response possible. When God awakens faith, moves us to repent, or enables an act of charity, he is drawing a real person into a real friendship. He enables our willing and acting, so that obedience can be our answer to his love. Prayer, sacramental life, and works of mercy are acts through which we make that answer. As we act, his grace continues to sustain us and to deepen the charity with which we respond.

A sober desire for heaven therefore includes the desire to be changed. We may long to be free from pain while guarding the resentment that makes another person's happiness unwelcome. We may wish to see God while refusing the truth he has already shown us about an injustice. These contradictions do not require us to abandon hope; they identify places where hope must become prayer and repentance. The heavenly good is God in his holiness. To desire him sincerely is gradually to consent to the purification of what prevents us from loving him and the people he gives us to love.

The discipline of receiving ordinary goods gratefully helps this consent. We can enjoy a meal without making comfort the measure of the day, accept praise without feeding upon it, and grieve an ending without declaring that all good has ended. In each case we receive a finite good without demanding that it secure our whole happiness. The freedom involved is sometimes costly because an attachment has become familiar and seems part of our identity. As those attachments loosen, we can discover how much of our conduct they had determined. The change can feel like loss even when it makes us freer to love.

